\section{Performance options}
\label{sec:performance}

Four settings decide most of the cost of a run, and only one of them
can change a number the user reads.  This chapter shows how to set
each one and what to expect.  The measured anchors behind the
recommendations are in the Manual (Chapter on performance and numerical
hygiene); the example script reproduces the routing and the two
implementation comparisons on a small system.

\subsection{Choosing the CI solver}

\texttt{canonicalize\_jax} and \texttt{canonicalize\_sa\_jax} take a
\texttt{ci\_solver} argument with three values:

\begin{itemize}
  \item \texttt{'auto'} (default): the solver is chosen from the
    determinant count $N_{\mathrm{det}} = \binom{2 n_{\mathrm{act}}}{n_{\mathrm{ele,act}}}$.
    At or above 2000 determinants the blocked solver is used, below it
    the plain solver.
  \item \texttt{'plain'}: the original inverse-iteration solver.  It is
    faster in eager execution on small spaces and returns a CI vector
    with a residual of order $10^{-8}$.
  \item \texttt{'blocked'}: the $S_z$-blocked solver.  It returns the
    exact eigenvector and works one block at a time, which is what
    makes large determinant spaces fit in memory.
\end{itemize}

\begin{lstlisting}[language=Python]
from diffmrc.core.nevpt2.impl import canonicalize_jax, resolve_ci_solver

# What would 'auto' select for this active space?
solver = resolve_ci_solver('auto', nact, nele_act)
print(solver)          # 'plain' below 2000 determinants, else 'blocked'

C_can, ci_can, eps = canonicalize_jax(
    C, h1e, eri, ci, ncore=ncore, nact=nact, nele_act=nele_act,
    ci_solver='auto')
\end{lstlisting}

Expected behaviour on the two measured anchors:

\begin{center}
\begin{tabular}{l r l p{0.34\textwidth}}
\toprule
Active space & $N_{\mathrm{det}}$ & \texttt{'auto'} selects & Why \\
\midrule
CAS(6,6) & 924 & \texttt{'plain'} & blocked is $1.8\times$ slower
  eager with no memory saving \\
CAS(8,8) & 12870 & \texttt{'blocked'} & peak memory
  $37.83 \to 6.78$~GB, and the plain jit graph does not compile \\
\bottomrule
\end{tabular}
\end{center}

The plain solver is not a smaller version of the blocked one: the
orbitals and orbital energies are bit-identical between the two, and
the CI vector differs only by a global phase and by the plain solver's
inverse-iteration residual.  A user who needs the exact pre-change CI
vector can pass \texttt{'plain'}; a user who forces \texttt{'blocked'}
on a small space pays a real cost, measured at $1.8$ times on 924
determinants and $36$ percent on the whole NEVPT2 theta chain at
CAS(2,2).

The flag affects only the internal re-solves of the canonicalization
routines.  The main CI solve of the production chain stays on the
plain solver, because substituting the blocked eigenvector there would
move the NEVPT2 correlation energy beyond the $10^{-12}$ acceptance
band.  In the state-averaged routine with more than one root the flag
has no effect, because no multi-root blocked solver exists.

\subsection{MO integral transforms}

Every MO integral transform in the differentiable chain goes through
\texttt{core.nevpt2.impl.\allowbreak stepwise\_mo\_ints}.  Use it in
place of the one-shot four-operand contraction:

\begin{lstlisting}[language=Python]
from diffmrc.core.nevpt2.impl import stepwise_mo_ints

# Use this:
h1e_mo, eri_mo = stepwise_mo_ints(C, h1e, eri)

# Not this:
# h1e_mo = jnp.einsum('ai,ab,bj->ij', C, h1e, C)
# eri_mo = jnp.einsum('pi,qj,rk,sl,pqrs->ijkl', C, C, C, C, eri)
\end{lstlisting}

The two forms are algebraically the same.  On the release tree the
stepwise form reproduces the H2O/6-31G CAS(2,2) energy and theta
gradient bit for bit, and the largest difference over the checked
systems is $5.7\times10^{-14}$~Ha on the linear CO case.  On a
96-orbital physical integral array it is also eleven to thirteen times
faster ($5.1$--$5.8$~s down to $0.42$--$0.47$~s).  On a small array
both forms finish in microseconds, so the gain is a production-size
effect rather than a small-system one.

The reason to keep the four steps is not only speed.  The one-shot
form was measured at an absolute accuracy of about $10^{-6}$ on the
large array, and a linear molecule's degenerate Fock pairs have gaps
near $10^{-15}$; the perturbation rotates the eigenvector basis inside
the degenerate subspace, which changed the theta gradient by eight
percent and made the Hessian-vector product NaN.

\subsection{jit with reference arrays as arguments}

The theta-only factories accept the reference arrays either as closure
constants or as runtime arguments.  Closure constants are compiled
into the traced graph, so a jitted wrapper must recompile for every
geometry.  Passing the arrays as arguments keeps one executable for
all geometries whose arrays share a shape and dtype.

The callable returned by \texttt{make\_nevpt2\_theta\_fn} exposes the
array-argument function as \texttt{component\_args} and the tuple to
pass as \texttt{arrays}:

\begin{lstlisting}[language=Python]
import jax

E_theta = make_nevpt2_theta_fn(mol, ncore, nact, nele, nmo, pairs,
                               scf_builder)
E_jit = jax.jit(E_theta.component_args)     # one compilation
E_ref = E_jit(theta, *E_theta.arrays)       # works for every geometry

# The orbital-point search accepts the same pair:
converge_orbital_point(E_theta.component_args, len(pairs),
                       arrays=E_theta.arrays)
\end{lstlisting}

On three geometries (the reference scaled by $-1$, $0$, and $+1$
percent) the closure form made three compilations and took $8.60$~s
for the three calls; the array form made one compilation and took
$2.67$~s.  At f-block window sizes a single compile is of order
fifteen minutes, which is where the difference matters.

The joint geometry-dependent entry \texttt{make\_nevpt2\_joint\_fn}
cannot be traced under \texttt{jax.jit} in this pyscfad build because
the SCF inside the call raises a tracer-conversion error.  The array
form is therefore for the fixed-geometry theta-only entry; the NEVPT2
correlation term is inside that jitted function and agrees with its
eager evaluation bit for bit.

\subsection{Running the example}

The script \texttt{examples/performance\_options.py} in the repository
demonstrates all four points on a small system.  From the repository
root:

\begin{lstlisting}[language=bash]
PYTHONPATH=src python examples/performance_options.py
\end{lstlisting}

It prints the routing decision for 924 and 12870 determinants, the
maximum difference between the stepwise and one-shot transforms on a
random $10^4$ array, the H2O/6-31G CAS(2,2) comparison of the three
solver settings, and the compile counts of the two jit forms.  The
recorded run ends with:

\begin{lstlisting}[language=]
CAS(6,6)  ndet=   924  auto -> plain
CAS(8,8)  ndet= 12870  auto -> blocked
max |d h1|  = 0.000e+00
max |d eri| = 1.332e-15
auto     vs plain: max|dC| = 0.0e+00  max|deps| = 0.0e+00  |dE_CAS| = 0.0e+00  max|dci| = 0.0e+00
blocked  vs plain: max|dC| = 0.0e+00  max|deps| = 0.0e+00  |dE_CAS| = 0.0e+00  max|dci| = 4.7e-10
array-argument form: jit objects = 1, compilations = 1
closure form:        jit objects = 3, compilations = 3
\end{lstlisting}

The H2O energy is $-61.314092285778507$~Eh in all three solver
settings; the blocked CI vector differs from the plain one by
$4.7\times10^{-10}$ after phase alignment, which is the plain solver's
inverse-iteration residual and not a solver error.

\subsection{What to record when timing}

Record the compilation count and the peak resident memory together
with any wall clock.  The first two are load-independent and decide
whether a run is feasible; wall clock on a shared node depends on the
load average and is only a relative quantity.  The example script
prints the compile counts from the jit objects' cache size, and the
Manual chapter records the peak memory of each route at the two
determinant anchors.
